%% Enroulement de l'axe réel sur le cercle trigonométrique (manuel p. 97).
%% Repère (O,I,J), cercle C vert de rayon 1 ; l'axe réel (rouge) est tangent au cercle en I ;
%% en l'enroulant, le réel 1 vient sur le point M et la longueur de l'arc IM vaut 1.
%% Sens direct (antihoraire, +) et sens indirect (horaire, en tirets).
\documentclass[tikz,border=4pt]{standalone}
\usepackage{liaisons}
\begin{document}
\begin{tikzpicture}[x=1.7cm,y=1.7cm,
  axe/.style={-{Stealth[length=5pt]}, line width=0.6pt},
  cercle/.style={solideE, line width=1.1pt},
  reel/.style={solideA, line width=1.1pt},
  sens/.style={-{Stealth[length=6pt]}, line width=0.8pt},
  rayon/.style={line width=0.6pt, black!70}]

%% ---------------- repère ---------------------------------------------
\draw[axe] (-1.95,0) -- (1.42,0);
\draw[axe] (0,-1.5) -- (0,1.5);
\node[font=\small, anchor=north east, inner sep=2pt] at (-0.03,-0.03) {$O$};

%% ---------------- cercle trigonométrique ------------------------------
\draw[cercle] (0,0) circle (1);
\node[font=\small, text=solideE, anchor=east, inner sep=3pt] at (200:1) {$\mathcal{C}$};

%% ---------------- l'axe réel, tangent en I ----------------------------
\draw[reel, -{Stealth[length=6pt]}] (1,-1.45) -- (1,2.3);
\foreach \yv/\lab in {-1/{$-1$}, 1/{$1$}, 2/{$2$}} {
  \draw[reel] (1,\yv) -- (1.12,\yv);
  \node[font=\small, text=solideA, anchor=west, inner sep=2pt] at (1.14,\yv) {\lab};}
\draw[reel] (1,0) -- (1.12,0);
\node[font=\small, text=solideA, anchor=west, inner sep=2pt] at (1.14,0) {$0$};

%% ---------------- points I, J et M ------------------------------------
\fill (1,0) circle (1.6pt);
\node[font=\small, anchor=north east, inner sep=2pt] at (0.99,-0.02) {$I$};
\fill (0,1) circle (1.6pt);
\node[font=\small, anchor=east, inner sep=3pt] at (0,1) {$J$};

%% M : point associé au réel 1 (angle de 1 radian), arc IM de longueur 1
\draw[rayon] (0,0) -- (57.2958:1);
\draw[solideA, line width=1.8pt] (0,0) ++(0:1) arc (0:57.2958:1);
\fill[solideA] (57.2958:1) circle (1.7pt);
\node[font=\small, anchor=south west, inner sep=1pt] at (57.2958:1.03) {$M$};
%% l'enroulement : la graduation 1 de l'axe réel vient sur M
\draw[-{Stealth[length=5pt]}, black!55, line width=0.6pt, dash pattern=on 2.5pt off 2pt]
     (0.96,1) to[bend right=25] (61:1.06);

%% ---------------- sens direct et sens indirect ------------------------
\draw[sens] (75:1.35) arc (75:130:1.35);
\node[font=\small, anchor=east, inner sep=2pt] at (-0.9,1.28) {$+$ sens direct};
\draw[sens, dash pattern=on 3pt off 2pt] (-75:1.35) arc (-75:-130:1.35);
\node[font=\small, anchor=east, inner sep=2pt] at (-0.9,-1.28) {sens indirect};
\end{tikzpicture}
\end{document}
